\documentclass[12pt]{article}
\usepackage{amsmath}
\usepackage{amsfonts}
\usepackage{amssymb}
\usepackage{geometry}
\usepackage{graphicx}
\geometry{a4paper, margin=0.8in}
\pagenumbering{gobble}

\title{Homework assignment for 02YMECH \\ Chapter: Rigid Body Motion}
\begin{document}
\date{}
\maketitle
\vspace{-0.8in}
\begin{enumerate}
    \item Calculate the moment of inertia of a uniform square plate of side length $b$ and mass $M$ about an axis that passes through its center of mass and is perpendicular to the plane of the plate.

    \item Consider a uniform solid cube of side length $b$ and mass $M$.
    \begin{enumerate}
        \item Compute the inertia tensor with respect to axes aligned with the edges of the cube, taking the origin of the coordinate system at the center of mass. Compare your result with the inertia tensor obtained in the lecture for a cube when the origin is chosen at one of its corners.
        \item Since the inertia tensor obtained in part (a) is diagonal in this coordinate system, under what conditions would the angular momentum vector be parallel to the angular velocity vector? Explain briefly.
        \item Interpret the physical meaning of the diagonal elements of the inertia tensor calculated in part (a). Is this result consistent with the moment of inertia of the square plate found in Question~1? Explain the connection.
        \item Suppose the cube is rotating freely in space with no external torques. About which axis (or axes) would you expect the rotation to occur? Based on this, discuss the significance of choosing the center of mass as the origin when calculating the inertia tensor.
    \end{enumerate}

    \item Consider a rigid body with principal moments of inertia
\[
I_1<I_2<I_3,
\]
and body-fixed principal axes
\((\mathbf e_1,\mathbf e_2,\mathbf e_3)\).
The motion is torque-free.

At time \(t=0\), the components of the angular velocity in the body frame
are
\[
\omega_1(0)=\omega_0,
\qquad
\omega_2(0)=\omega_0,
\qquad
\omega_3(0)=0,
\]
with \(\omega_0>0\).

\begin{enumerate}
  \item[(a)] Write Euler’s equations for this system and show that both the
        kinetic energy \(T\) and the squared angular momentum \(L^2\) are
        conserved.

  \item[(b)] Using only these two conserved quantities, eliminate
        \(\omega_2\) and \(\omega_3\) to obtain a first-order differential
        equation for \(\omega_1(t)\) of the form
        \[
        \dot{\omega}_1^{\,2}=f(\omega_1),
        \]
        and determine the explicit polynomial \(f\).

  \item[(c)] Determine the range of allowed values of \(\omega_1(t)\) and show
        that the motion of \(\omega_1\) is periodic.

  \item[(d)] Express the period of \(\omega_1(t)\) as an integral,
        clearly stating the change of variables used. Do not evaluate the
        integral.
\end{enumerate}
\end{enumerate}


\end{document}
